% Continuation of D.2.48: the full invariant, with both exceptional moduli.
\begin{proof}[Proof of the complete invariant and classification]
\emph{A global finite-coefficient class, retaining boundary torsion.}
At weight one the regular bundle $V_1$ over the peripheral
graph has free $H_1$ and $H_2$. Its degree-two group has the
primitive fibre quotient generated by $s$, together with the
five marked graph generators in (D.WE11). Choose the integral
cohomology class $\alpha_1$ which evaluates as $\xi_1$ on the
fibre part and zero on every graph generator.

On the weight-one cusp use the actual toric line with support
\[
                 \psi(r,s)=-\frac{r(r-1)}2+3s(s-1).
 \tag{D.WI17}
\]
Its multiplier windings give $\xi_1$ in the positive period
marking used in (D.WI7). Its boundary value agrees integrally
with $\alpha_1$: their fibre restrictions agree, and both
vanish on the two angular Wang sweeps, which bound in the
cusp. These are all the remaining generators of $H_2(M_1)$.
Also $H_1(M_1)$ is free, so the integral universal coefficient
theorem leaves no hidden $\operatorname{Ext}$ difference.

Before the finite caps are attached, the core and cusp have
the genuine degree-$d$ period covering defined by
$\gamma_d=d\gamma_1$. Pull back both classes along that
cover. They then agree on the entire boundary $M_d$, not
only on its regular-fibre restriction. In particular this
argument retains the $\mathbb Z/d$ term in $H_1(M_d)$ and
the resulting torsion in $H^2(M_d)$. The core pullback
$\alpha_d$ has fibre functional $\xi_d$ and zero values on
the five graph generators: the covering fixes the first
three such generators and multiplies the two $u_i$ generators
by $d$.

At a finite cap let $\pi:Y_i\longrightarrow Q_i$ be its
actual normal circle bundle. The obstruction for
$\alpha_d\bmod d$ to extend across the disk bundle is its
Gysin image
\[
                  \pi_!(\alpha_d\bmod d)\in
                         H^1(Q_i;\mathbb Z/d).
 \tag{D.WI18}
\]
Its values on the complete marked generators
$a,\gamma,g$ of $H_1(Q_i)$ are precisely the first-row
coefficients of (D.WE3):
\[
 \begin{array}{c|ccc}
       &a&\gamma&g\\ \hline
 i=0&-2dn_0&-6dr_0&0\\
 i=1&-3dn_1&-6dr_1&0.
 \end{array}
 \tag{D.WI19}
\]
The root-loop column $g$ is retained. This is the Thom
definition of fibre integration: a loop in $Q_i$ sweeps its
normal circle, and pairing with $\alpha_d$ gives the displayed
actual normal-column coefficient. Every coefficient is zero
modulo $d$. Since $H^1(Q_i;\mathbb Z/d)$ is
$\operatorname{Hom}(H_1(Q_i),\mathbb Z/d)$, (D.WI18) vanishes,
even when the integral normal line has torsion. The circle
Gysin sequence therefore supplies an extension over each
finite cap. Mayer--Vietoris now glues them with the pulled-back
cusp line to a class $u\in H^2(X;\mathbb Z/d)$.
Only the core and cusp were covered; neither finite fill was
assumed to be a cover of a weight-one fill.

On every normalized cusp component the integral representative
of its restriction is
\[
              \eta=-5H+6E_1-E_3,\qquad
       \eta(E_1,E_2,E_3)=(-6,0,1),\qquad
       \eta(f)=1,\qquad \eta^2=-12.
 \tag{D.WI20}
\]
The three degrees follow directly by taking the differences
of (D.WI17) across the internal, diagonal, and scalar-fixed
walls, as in D.1.9. In particular $u(s)=1$.
Choose free lifts in $H_2(X)$ annihilated by $u$, subtracting
appropriate multiples of $s$. This specifies the splitting
in (D.WI14)--(D.WI16). The integral classes $x_i$ form its
dual free basis, and $u$ is the standard cyclic coordinate.

\emph{The full mixed coordinates.}
The first mixed value in (D.WI15) follows from (D.WI3) and
$\langle u\kappa,[X]\rangle=1$. The Pontryagin square of the
reduction of an integral degree-two class is its ordinary
square with the full target modulus $2d$; this is the
$p=2$ case of the actual cochain construction
\cite[equations (2.1), (2.8), and (3.4), pp.267--268]{ffThomas56}.
By naturality and the projection formula for the normalized
component fundamental class,
\[
 \langle \mathcal Q_2(u)x_i,[X]\rangle
       =\langle\mathcal Q_2(j_i^*u),[dP_6]\rangle
       =\eta^2=-12\pmod {D_2}.
 \tag{D.WI21}
\]
This computes the coefficient modulo $2d$ when $d$ is even;
reducing the ordinary square merely modulo $d$ would lose
precisely the extra order-two information. The relation
$\sum_i x_i=0$ is consistent with (D.WI21), since
$12d$ is divisible by $D_2$.

\emph{The pure coordinate.}
The following actual scalar-circle calculation evaluates
the full Pontryagin third power. It uses the class just
constructed and its specified component restrictions.
% Primary-operation input for maps-owned D.2.48; no separate theorem number.
\subsubsection*{The full Pontryagin cube of the weighted Engel torsion class}
We use the actual class $u\in H^2(X_d;\mathbb Z/d)$ whose
restrictions to the normalized cusp components are specified
in D.2.48. In particular its degree on every scalar-fixed
rational curve $C_j$, $j\in\mathbb Z/d$, is one, and its
component restriction is represented by the integral toric
support
\[
                  \psi(r,s)=-\frac{r(r-1)}2+3s(s-1).
                                                               \tag{D.TP1}
\]
The calculation below uses these actual restrictions. Merely
choosing an arbitrary degree-two class that is a unit on
torsion would not fix its pure coordinate in the presence
of the free summand.

The scalar circle action is semifree. On a finite cap of
order $m_i$ its gamma translation has primitive numerator
$n_i$, with $\gcd(m_i,n_i)=1$. A nontrivial root power cannot
be compensated by the scalar direction, because that would
require $j n_i/m_i$ integral. Thus these caps are scalar-free.
The cusp has exactly the $d$ fixed curves $C_j$; their normal
bundles are $\mathcal O(-1)\oplus\mathcal O(-1)$ with weights
$(1,-1)$, and the other cusp orbits are free. These are the
marked geometric inputs of the weighted family.

Because $X_d$ is simply connected, $u$ has an equivariant
lift $\widehat u\in H^2_{S^1}(X_d;\mathbb Z/d)$.
In the Borel spectral sequence the only possible
differentials out of ordinary degree two have targets
involving $H^1(X_d;\mathbb Z/d)$ or $H^3(BS^1;\mathbb Z/d)$,
and both vanish. Write
\[
         \widehat u|_{C_j}=x+\lambda_j t,
         \qquad x=[C_j]^*,\quad t\in H^2(BS^1;\mathbb Z/d).
                                                               \tag{D.TP2}
\]
The difference of any two equivariant lifts of a fixed
ordinary class on the simply connected normalized
$dP_6$ is a multiple of $t$. Hence the differences of
the characters on its two scalar-fixed curves can be
read from any toric equivariant lift of (D.TP1).
The fixed edges are the translates of the edge from
$(0,j)$ to $(0,j+1)$; the support slope in their scalar
direction is
\[
               \psi(0,j+1)-\psi(0,j)=6j.
                                                               \tag{D.TP3}
\]
The fibre character of the associated monomial
trivialization is the negative of this slope, up to
one global character convention. Equivalently the
intervening weight-one orbit curve has degree $-6$.
Thus, for one common sign $\epsilon=\pm1$,
\[
                    \lambda_j=\lambda_0+6\epsilon j
                                 \quad\text{in }\mathbb Z/d.
                                                               \tag{D.TP4}
\]
This equation is compatible with the period $j\mapsto j+d$.
It does not assert that the integral support line itself
descends with an integral global circle linearization:
the needed character equality is modulo $d$ and follows
from the existing global equivariant cohomology lift.

Suppose now that $3\mid d$. Thomas's operation is defined
with precisely the required full coefficients:
\[
 \mathcal P_3:H^2(-;\mathbb Z/d)
                    \longrightarrow H^6(-;\mathbb Z/(3d)).
                                                               \tag{D.TP5}
\]
This is \cite[equation (2.1), p.267]{ffThomas56} with
$p=3$ and $m=d/3$. On the reduction of an integral
class $z$, it equals the reduction of $z^3$ modulo
$3d$. Indeed the correction term in Thomas's actual
cochain formula (3.4) is zero for an integral cocycle,
and its first term is the ordinary threefold diagonal
product \cite[equations (3.4)--(3.6), p.268]{ffThomas56}.
Each restriction (D.TP2) has such an integral lift.

Apply equivariant localization with coefficients
$\mathbb Z/(3d)$ to $\mathcal P_3(\widehat u)$.
No field assumption is needed here: the complement
is scalar-free, and each normal Euler class becomes
a unit after inverting $t$. On a fixed curve it is
\[
                    (-x+t)(-x-t)=-t^2,
\]
because $x^2=0$. Integrating the term of ordinary
curve degree two gives
\[
 \left\langle\mathcal P_3(u),[X_d]\right\rangle
                      =-3\sum_{j=0}^{d-1}\lambda_j^2
                                  \quad\text{in }\mathbb Z/(3d).
                                                               \tag{D.TP6}
\]
The right side is independent of the integer
representatives of the characters: $3\mid d$ makes
the cube of an integral lift well defined modulo
$3d$. The base cohomology polynomial ring injects
into its localization at $t$, so the degree-zero
identity is an identity in the original coefficient
ring, not just a localized statement.

For arbitrary integers $d>0,\lambda_0$ and either
sign in (D.TP4), the exact identity is
\[
 \sum_{j=0}^{d-1}(\lambda_0+6\epsilon j)^2
   =d\left(\lambda_0^2+6\epsilon\lambda_0(d-1)
                         +6(d-1)(2d-1)\right).              \tag{D.TP7}
\]
Substitution into (D.TP6) proves
\[
                 \left\langle\mathcal P_3(u),[X_d]\right\rangle=0
                                  \quad\text{in }\mathbb Z/(3d).
                                                               \tag{D.TP8}
\]
Thus the complete pure cyclic fundamental-class coordinate
is zero, including the order-three subgroup invisible to
the ordinary mod-$d$ cube. For $3\nmid d$ the same
localization calculation with the ordinary cube gives
zero modulo $d$. Mixed coordinates are separate from
this calculation; in particular it does not discard
the order-$2d$ mixed Pontryagin-square terms when $d$
is even.

The universal integer identities, support slope,
single-curve Euler/cube calculation and an explicitly
bounded supplementary check are executed in
\path{code/fourfolds_D48_thomas.py}.
The supplementary range is $d=3,6,\ldots,120$,
$\lambda_0=0,\ldots,d-1$ and both signs, giving
$4920$ exact checks. The arbitrary-$d$ conclusion
comes from (D.TP7), not that finite range.


\emph{Why these detectors are complete.}
We include the low-degree integral target calculation to
exclude an undetected finite subgroup. For a cyclic group,
\[
 \begin{aligned}
 H_2(K(\mathbb Z/d,2);\mathbb Z)&=\mathbb Z/d,\\
 H_4(K(\mathbb Z/d,2);\mathbb Z)&=\mathbb Z/D_2,\\
 H_6(K(\mathbb Z/d,2);\mathbb Z)&=\mathbb Z/D_3 .
 \end{aligned}
 \tag{D.WI22}
\]
The degree-four assertion follows from Whitehead's exact
sequence: on $K(A,2)$ its two adjacent homotopy groups
vanish, so $H_4=\Gamma(A)$
\cite[Theorem 1 and Section 2]{mapsWhitehead50}.
For $A=\mathbb Z/d$, the universal quadratic generator
$q(1)$ has relations $d^2q(1)=2dq(1)=0$.
Conversely $q(j)=j^2$ modulo $\gcd(d^2,2d)=D_2$
is well defined and quadratic. Thus these are its complete
relations and the image of $\mathbb {CP}^2$ generates.
The ordinary or Pontryagin square evaluates on that image
as one modulo $D_2$, so it detects the entire group.

For degree six, Zhubr's epimorphism in Section 3.11 and
his two free bordism generators in Section 3.12 show that
the spin-realizable part for a cyclic coefficient group
is generated in homology by the image of
$[\mathbb {CP}^3]$: his other generator has zero
fundamental-class image. For odd-primary cyclic groups
this is all of $H_6$, since the only obstruction in his
Lemma 3.7 has mod-two target. For a two-power cyclic
group Section 3.12 explicitly gives
$\ker\alpha=H_6=\mathbb Z/2^a$
\cite[Lemmas 3.7 and 3.11, and Section 3.12,
pp.805--807]{manZhubr75}.
The decomposition of $\mathbb Z/d$ into its relatively
prime primary factors gives the corresponding product
of Eilenberg--MacLane spaces. All cross tensor and
$\operatorname{Tor}$ terms vanish, since their positive
homology has relatively prime orders. Hence $H_6$ is
cyclic and is generated by the image $z$ of
$[\mathbb {CP}^3]$ for every $d$.

The order of $z$ divides both $d^3$ and $3d$.
Indeed the class $d h$ on $\mathbb {CP}^3$ becomes zero
modulo $d$, while its integral cube is $d^3$.
On $\mathbb {CP}^1\times\mathbb {CP}^2$ the class
$d x+y$ reduces to a class factoring through
$\mathbb {CP}^2$, whose degree-six image is zero,
while its integral cube evaluates to $3d$.
Thus the order divides
$\gcd(d^3,3d)=d\gcd(d,3)=D_3$.
The opposite bound is detected by the ordinary cube when
$3\nmid d$, or by Thomas's operation when $3\mid d$:
on the reduced hyperplane class of $\mathbb {CP}^3$
the full target value is one
\cite[equation (2.8), p.267]{ffThomas56}.
This proves the last assertion of (D.WI22) with its
stated perfect detector.

Finally
$K(\mathbb Z^r\oplus\mathbb Z/d,2)
 =(\mathbb {CP}^{\infty})^r\times K(\mathbb Z/d,2)$.
The first factor has torsion-free homology in even
degrees and zero homology in odd degrees. The integral
Kunneth formula has no $\operatorname{Tor}$ term and
gives exactly (D.WI14). Integral cubic monomials,
$u x_i x_j$, $\mathcal Q_2(u)x_i$, and
$\mathcal Q_3(u)$ are the respective perfect coordinate
detectors. Therefore (D.WI15) computes all of $\mu$,
including the extra factors two and three.

\emph{Application of smooth classification.}
The manifold is simply connected and spin, has
$H_2=\mathbb Z^{d-1}\oplus\mathbb Z/d$ and middle Betti
number zero by D.2.34, refined class zero by the first
proof, and the complete fundamental-class invariant
(D.WI15). These are exactly the four invariants in
Zhubr's Theorem 1, with the definitions on pp.794--795
\cite[Theorem 1]{manZhubr75}. Their values are independent
of every marked parameter in (D.WI1).
Existence is supplied by the constructed manifolds,
and the injective classification gives one oriented
diffeomorphism type for each $d$. This proves (D.WI16)
by an applied complete classification, rather than by
ordinary cohomology or an asserted connected-sum model.
\end{proof}
